% !TEX root = ../../../tree_notes.tex
% =====================================================================
%  13.4  Distance-indexed DP: histograms over depth in the subtree
%        sec:c10-distance
%  Exports: def:c10-dist-hist  thm:c10-depth-xor  lst:c10-depth-hist
%           fig:c10-distk  ex:c10-k2 (Ch.11's forward contract)
%           prob:c10-distk-solved (CF 570D)  lst:c10-570d-soln
%  Discharges Ch.16's forward-contract references sec:c10-distance (from
%  ch11/04_paths_lca.tex L21) AND ex:c10-k2 (from ch11/01_ancestor_lca.tex
%  L76 -- the "degenerate rows ride on l itself, one arm of length 0"
%  census).  Solves CF 570D with a per-depth XOR accumulator maintained
%  by a backtracking DFS; knapsack-budget analysis previewed, deferred
%  to Ch.14 (c27).  All numbers on T hand-verified.
% =====================================================================
\section{Distance-Indexed DP: Histograms over Depth in the Subtree}\label{sec:c10-distance}

\begin{previouslybox}
\begin{enumerate}[nosep, leftmargin=1.6em]
  \item Chapter~12's \secref{c09-measures} computes $\depth(v)$ from the
        parent as a top-down parameter and $\height(v),\,\mathit{sz}(v)$
        from the children as a bottom-up fold. If you keep an
        \emph{array indexed by depth} while the DFS runs, what does the
        array hold at the moment the recursion is at $v$? (aggregate over
        nodes already visited at that depth -- the shape of the state
        this section makes precise)
  \item \thmref{c10-pairtop-count} harvested pairs at the peak with
        $\bigO{k}$-per-child histogram merges. When the query asks
        ``the aggregate at depth $d$ in $v$'s subtree'' for many
        $(v,d)$ pairs, is the harvest-at-peak machinery still the right
        tool, or does a different DP state do better? (a different state
        -- per-depth \emph{accumulator} on the DFS path)
  \item Chapter~11's \exref{c11-census} showed that
        $\lca(u,v)=u$ when $u \preceq v$ -- a ``degenerate'' row of the
        LCA table. \secref{c10-harvest}'s peak census has an analogous
        row; name it and state its contribution to the count. (the arm
        of length~$0$ at the peak, i.e.\ one endpoint \emph{is} $\ell$;
        contributes $\sum_{c}H_{c}[k-1]$ to the distance-$k$ total)
  \item CF~161D's answer for $k=2$ on $T$ is $10$, with $7$ of those
        pairs harvested at a peak that equals one of their endpoints.
        Where in the histogram recurrence do these $7$ live? (in the
        row $\mathit{H}[u][0]=1$ set \emph{before} any child merge)
\end{enumerate}
\end{previouslybox}

\begin{topicintro}
Section~13.1 folded a \emph{constant-size} tuple per node; 13.2 doubled it;
13.3 indexed the state by the \emph{length} of an arm and produced an
$\bigO{k}$ histogram. This section makes that index the star: the DP state
at a node is a \emph{histogram over depths in the subtree}, and the merge
across children is a knapsack. The \propref{c10-pairtop} harvest identity
stays valid and is the reason the merge is legal; the new material is how
to answer \emph{per-node, per-distance} queries with a single DFS and how
the histogram merge's cost trades against tree height.

Two solved shapes populate this section. The first is the
\emph{distance-$k$ census} of \secref{c10-harvest}, re-derived here in
histogram language and traced as \exref{c10-k2} on the running tree $T$
-- the census that Chapter~11's \exref{c11-census} ancestor/non-ancestor
distinction leans on when it says ``the degenerate rows ride on $\ell$
itself, one arm of length~$0$''. That example discharges
\exref{c10-k2}. The second is CF~570D (\probref{c10-distk-solved}),
\emph{Tree Requests}: per-node, per-depth character queries, answered in
$\bigO{n+m}$ with a \textbf{backtracking DFS} and a per-depth XOR mask --
a cleaner variant than the histogram merge because queries are offline.
Together they bracket the design space.

\textbf{Prerequisites.} \secref{c10-harvest}'s \propref{c10-pairtop} and
the arm histogram; \secref{c10-subtrees}'s fold template; Chapter~12's
\secref{c09-measures} for the $\depth/\height/\mathit{sz}$ tables on $T$;
\emph{not} assumed: Euler-tour techniques (Ch.24's \chref{c18} owns the
tin/tout calculus) and DSU-on-tree small-to-large merges (Ch.17's
\chref{c28} owns the $\bigO{n\log n}$ budget). Both appear as forward
pointers here.

\textbf{After this section you can}
\begin{itemize}[nosep]
  \item state a per-node histogram $\mathit{H}_v[\cdot]$ over depths and
        write the child-merge that shifts $+1$ on the way up; \hfill\lv{3}
  \item answer offline per-node, per-depth queries in $\bigO{n+m}$ via a
        backtracking DFS with a path-side accumulator; \hfill\lv{4}
  \item name the ``arm of length~$0$'' degenerate row in the $k$-census
        and cite \exref{c10-k2} when Ch.16's LCA machinery relies on it;
        \hfill\lv{3}
  \item bound the knapsack merge budget $\sum_v \deg(v)\cdot k$ and
        recognise when it exceeds the constraint. \hfill\lv{4}
\end{itemize}
\end{topicintro}

\subsection{Intuition}

A \emph{scalar per node} is enough for aggregates that fold to the root
(\secref{c10-subtrees}); a \emph{pair per node} for take/leave decisions
(\secref{c10-twostate}); an \emph{array per node} once the query carries
an extra integer parameter. The array's index is the parameter. Concretely:

\begin{itemize}[nosep]
  \item CF~161D asks ``pairs with $d(u,v)=k$'' -- index by arm length $j$
        from the peak, and the answer lives at $j + j' = k$ (\thmref
        {c10-pairtop-count}, already harvested in \secref{c10-harvest}).
  \item CF~570D asks ``at depth $h$ in subtree $v$, are the letters a
        palindrome?'' -- index by absolute depth $d$ from the root, keep
        a per-depth XOR mask, harvest at $v$'s subtree boundaries.
  \item A general ``sum over pairs with $F(u,v)$ additive in $d(u,v)$''
        uses the peak histogram of \thmref{c10-pairtop-count} with a
        different $\mathrm{combine}$.
\end{itemize}

Two algorithmic shapes recur. The \term{bottom-up histogram merge}
builds $\mathit{H}_v[j]$ for every $v$ explicitly by folding children's
$\mathit{H}_c[j-1]$ into $\mathit{H}_v[j]$; this is CF~161D's skeleton
and the ``harvest at the peak'' of \secref{c10-harvest}. Its total cost
is $\sum_v \deg(v)\cdot|\mathit{H}_v|$, which is $\bigO{nk}$ for a
bounded index range $k$ but degrades to $\bigO{n\cdot\height}$ when the
histogram must run up to $\height(v)$. The \term{backtracking-DFS
accumulator} keeps a \emph{single} global array $\mathit{X}[d]$ indexed
by depth and updates it on entry/exit of each node; queries attached to
$v$ are answered at entry or exit time. Its total cost is $\bigO{n+m}$
for $m$ offline queries. Both shapes are the same DP -- a state per
depth -- paid for differently.

\subsection{Definitions and Notation}

\begin{definition}{Depth-indexed subtree state (histogram)}{c10-dist-hist}
Fix a root $r$ and let $\depth(\cdot)$ be the edge-count depth from $r$.
A \term{depth-indexed subtree state} at $v$ is a map
$\mathit{H}_v\colon \{0,1,\dots,D_v\}\to \mathfrak{A}$ into some
\emph{aggregated} set $\mathfrak{A}$ (an integer count, a bitmask, a
small struct) satisfying
\[
  \mathit{H}_v[j]
  \;=\;
  \Phi\!\Big(\,v\text{'s own contribution at }j,\;\;\big\{\mathit{H}_c[j-1]:c\in\children(v)\big\}\,\Big),
\]
where $\Phi$ is a fixed combiner ($+$ for counts, $\oplus$ for XOR
masks, $\max$ for extrema). Here $D_v$ is the histogram range we retain;
$D_v$ may be $\le \height(v)$ (full), $\le k$ (bounded), or a global
$\le D$ (used when the state is a \emph{path}-side accumulator rather
than a per-node table).
\end{definition}

\noindent\textbf{In plain words.} The state at $v$ is an \emph{array},
indexed by depth, whose entry at $j$ is the fold of $v$'s subtree at
distance $j$ below $v$. The recurrence shifts by $+1$ at every edge
because children sit one level down.

\begin{theorem}{Offline per-node, per-depth queries via a prefix accumulator}{c10-depth-xor}
Let $Q$ be a set of $m$ offline queries of the form ``read the fold
$\bigoplus_{u \in \mathrm{subtree}(v),\,\depth(u)=h} \mathit{bit}(u)$'',
each tagged with a node $v$ and depth $h$. If the combiner $\Phi$ is a
\emph{group operation} on $\mathfrak A$ (invertible), Algorithm
\coderef{c10-depth-hist} answers all $m$ queries in $\bigO{n + m \cdot |\mathfrak{A}|}$
time and $\bigO{n + m}$ space during a single DFS.
\end{theorem}
\begin{proofbox}
Maintain a global array $\mathit{X}[d]$ that accumulates $\Phi$-applications
of every node visited in preorder at depth $d$, and is \emph{never}
restored. Fix a query $(v,h)$. Take the value $\mathit{X}[h]_{\mathrm{in}}$
as of $v$'s entry, before any of $v$'s subtree has been processed; the
value $\mathit{X}[h]_{\mathrm{out}}$ at $v$'s exit equals
$\mathit{X}[h]_{\mathrm{in}} \mathbin{\Phi} \bigoplus_{u \in \mathrm{subtree}(v),\,\depth(u)=h} \mathit{bit}(u)$,
because the DFS visits every node in $v$'s subtree (and only those)
between the two snapshots and only nodes at depth $h$ mutate
$\mathit{X}[h]$. Applying the group inverse gives the aggregate:
$\mathit{X}[h]_{\mathrm{out}} \mathbin{\Phi^{-1}} \mathit{X}[h]_{\mathrm{in}}$. Total cost:
$\bigO{1}$ enter per node plus $\bigO{1}$ inverse-and-read per query. \qed
\end{proofbox}

\begin{insightbox}
The theorem's requirement that $\Phi$ be invertible is what lets the
accumulator \emph{subtract} the pre-subtree prefix from the post-subtree
prefix. For CF~570D, $\Phi=\oplus$ on 26-bit masks and $\oplus$ is its
own inverse -- trivial. For non-invertible combiners ($\max$,
count-of-max) the difference $\mathit{X}_{\mathrm{out}} - \mathit{X}_{\mathrm{in}}$
is not defined and Shape~B collapses; one falls back to the
\emph{bottom-up histogram merge} of \thmref{c10-pairtop-count}, at
$\bigO{nk}$ cost or $\bigO{n\log n}$ with the small-to-large refinement
Ch.17's \chref{c28} owns.
\end{insightbox}

\subsection{The Algorithm: Two Shapes}

\paragraph{Shape A: bottom-up histogram merge.}
This is \coderef{c10-peak} from \secref{c10-harvest} re-stated in
histogram language:

\begin{lstlisting}[caption={Bottom-up depth histogram merge (CF 161D shape): each child merged after harvest}, label={lst:c10-distk-hist}]
// adj[u] = neighbours; call hist(r, -1, K). Writes ans.
vector<vector<int>> adj;
vector<vector<int>> H;                     // H[u][j] = count of nodes at dist j in u's subtree
long long ans; int K;

void hist(int u, int p) {
    H[u].assign(K + 1, 0);
    H[u][0] = 1;                           // u itself at dist 0 -- the "arm of length 0" seed
    for (int v : adj[u]) if (v != p) {
        hist(v, u);                        // child histogram final first
        // harvest pairs whose peak is exactly u, one endpoint on
        // the already-merged side and the other in v's subtree
        for (int i = 0; i <= K - 1; ++i)
            ans += (long long)H[u][i] * H[v][K - 1 - i];
        // shift-merge v into u (each node in v's subtree gains 1 edge)
        for (int j = K; j >= 1; --j)
            H[u][j] += H[v][j - 1];
    }
}
\end{lstlisting}

\noindent Cost: $\bigO{nk}$. Every node owns $\le K+1$ histogram slots;
every edge triggers $2K$ arithmetic ops. This is \thmref
{c10-pairtop-count} rewritten as a merge -- the ``harvest then shift''
discipline is identical, so pairs at the peak are still never
double-counted and the length-$0$ arm at $u$ is still captured by
$\mathit{H}[u][0]=1$ before the child loop.

\paragraph{Shape B: backtracking-DFS accumulator (CF 570D).}

\begin{lstlisting}[caption={Per-depth prefix XOR along the DFS preorder; each query takes an entry snapshot and reads the exit XOR-difference}, label={lst:c10-depth-hist}]
// adj[u] = children; s[u] = character (0..25); depth[u] from root (root=1)
// X[d] = running 26-bit XOR mask over every node visited in DFS preorder
//        at depth d so far. NEVER restored -- this is the whole trick.
vector<vector<int>> adj;
vector<int> s, depth, X;
vector<vector<pair<int,int>>> query_at;   // query_at[v] = list of (h, qid)
vector<int> in_val;                        // per-query entry snapshot of X[h]
vector<char> ans;

void dfs(int u) {
    for (auto& [h, qid] : query_at[u]) in_val[qid] = X[h];    // snapshot
    X[depth[u]] ^= (1 << s[u]);                                // add u
    for (int v : adj[u]) { depth[v] = depth[u] + 1; dfs(v); }  // recurse
    for (auto& [h, qid] : query_at[u]) {
        int mask = X[h] ^ in_val[qid];                         // subtree XOR
        ans[qid] = (__builtin_popcount((unsigned)mask) <= 1);
    }
    // NO restore -- u's bit stays in X for the parent's aggregate.
}
// caller: depth[1] = 1; X.assign(n+2, 0); in_val.resize(m); dfs(1);
\end{lstlisting}

\noindent Cost: $\bigO{n + m}$, one $\bigO{1}$ XOR per node and
one $\bigO{1}$ popcount per query. This is \thmref{c10-depth-xor} with
$\Phi=\oplus$. Note the \emph{asymmetry} with Shape~A: there is no
corresponding ``leave'' line. Every node's contribution stays in $\mathit{X}$
for the whole rest of the DFS; that is exactly what makes the difference
$\mathit{X}[h]_{\mathrm{out}} \oplus \mathit{X}[h]_{\mathrm{in}}$ equal the subtree aggregate.

\begin{complexitybox}
\textbf{Shape A} is $\bigO{nk}$ time, $\bigO{nk}$ memory. Use it when the
index is a \emph{pair-level} parameter (``count pairs at distance $k$''
-- CF~161D). \textbf{Shape B} is $\bigO{n+m}$ time, $\bigO{n}$ memory.
Use it when queries are \emph{node-level} and $\Phi$ is invertible (CF
570D). When the index range is unbounded and $\Phi$ is not invertible,
neither shape fits: escalate to DSU-on-tree (Ch.17's \chref{c28}) or to
centroid decomposition (Ch.24's \chref{c18}). \thmref{c10-depth-xor}'s
group-invertibility clause is the honest reason CF~570D is easier than
it looks -- an equivalent statement with $\max$ instead of XOR is a
different problem.
\end{complexitybox}

\subsection{Worked Example: The $k=2$ Census on $T$, Grouped by Peak}\label{subsec:c10-k2-census}

\begin{example}{The $k=2$ distance-pair census on $T$, degenerate rows flagged}{c10-k2}
Take $T$ from \figref{c10-runtree-dp}, root at~$1$. Its edge list is
$1\!-\!2,\,1\!-\!3,\,2\!-\!4,\,2\!-\!5,\,3\!-\!6,\,3\!-\!7,\,5\!-\!8,\,
3\!-\!6,\,7\!-\!9,\,9\!-\!10$ (Ch.12's \secref{c09-measures} table on $T$
reused verbatim). Enumerate every unordered pair $\{u,w\}$ with $d(u,w)=2$
and file each under $\ell=\mathrm{peak}(u,w)=\mathrm{LCA}(u,w)$. Mark a
row \emph{degenerate} if one arm has length~$0$ (i.e., one endpoint IS
the peak -- the shortcut case Ch.16's \exref{c11-census} ancestor/non-
ancestor split relies on). The census is \ref{tab:c10-k2-table}; the
running total is $\mathbf{10}$, matching \exref{c10-harvest-run}'s
histogram sum.
\end{example}

\begin{table}[H]
\centering
\footnotesize
\renewcommand{\arraystretch}{1.2}
\begin{tabular}{@{}c l l c c@{}}
\toprule
$\ell$ & pair $\{u,w\}$ & arms at $\ell$ & $\mathit{len}(a)+\mathit{len}(b)$ & degenerate? \\ \midrule
$1$ & $\{1,4\}$ & $(0)+(1\!\to\!2\!\to\!4=2)$ & $0+2=2$ & \textbf{yes} \\
$1$ & $\{1,5\}$ & $(0)+(1\!\to\!2\!\to\!5=2)$ & $0+2=2$ & \textbf{yes} \\
$1$ & $\{1,6\}$ & $(0)+(1\!\to\!3\!\to\!6=2)$ & $0+2=2$ & \textbf{yes} \\
$1$ & $\{1,7\}$ & $(0)+(1\!\to\!3\!\to\!7=2)$ & $0+2=2$ & \textbf{yes} \\
$1$ & $\{2,3\}$ & $(1\!\to\!2=1)+(1\!\to\!3=1)$ & $1+1=2$ & no (cross-child) \\ \midrule
$2$ & $\{2,8\}$ & $(0)+(2\!\to\!5\!\to\!8=2)$ & $0+2=2$ & \textbf{yes} \\
$2$ & $\{4,5\}$ & $(2\!\to\!4=1)+(2\!\to\!5=1)$ & $1+1=2$ & no (cross-child) \\ \midrule
$3$ & $\{3,9\}$ & $(0)+(3\!\to\!7\!\to\!9=2)$ & $0+2=2$ & \textbf{yes} \\
$3$ & $\{6,7\}$ & $(3\!\to\!6=1)+(3\!\to\!7=1)$ & $1+1=2$ & no (cross-child) \\ \midrule
$7$ & $\{7,10\}$ & $(0)+(7\!\to\!9\!\to\!10=2)$ & $0+2=2$ & \textbf{yes} \\ \bottomrule
\end{tabular}
\captionsetup{type=table}
\caption{The $k=2$ census on $T$: $7$ degenerate rows riding on $\ell$
itself, $3$ strict cross-child rows, total $10$. Peaks $4,5,6,8,9,10$
carry no $k=2$ row: $5$ and $7$ each have a single child, so no
cross-child pair at distance~$2$ fits; the leaves have no depth-$2$
descendant at all. Every row with a bold ``yes'' is one the LCA
\emph{ancestor-of} case licenses (Ch.16's \propref{c11-chain}: if
$u\preceq v$ then $\lca(u,v)=u$).}\label{tab:c10-k2-table}
\end{table}

\begin{figure}[H]
\centering
\begin{tikzpicture}[diagram, scale=0.88, transform shape,
    tnode/.style={circle, draw=primaryblue, line width=0.9pt, fill=primaryblue!8,
                  minimum size=6.2mm, font=\small\bfseries, text=primaryblue!60!black},
    peak/.style={circle, draw=accentorange, line width=1.1pt, fill=accentorange!25,
                 minimum size=6.6mm, font=\small\bfseries, text=accentorange!70!black},
    arm0/.style={draw=accentorange, very thick},
    armX/.style={draw=successgreen!70!black, very thick}]
  \node[peak] (n1) at (0,0) {$1$};
  \node[tnode] (n2) at (-3,-1.6) {$2$};
  \node[peak] (n3) at ( 3,-1.6) {$3$};
  \node[tnode] (n4) at (-5.2,-3.2) {$4$};
  \node[tnode] (n5) at (-1.6,-3.2) {$5$};
  \node[tnode] (n6) at ( 1.6,-3.2) {$6$};
  \node[peak] (n7) at ( 5.2,-3.2) {$7$};
  \node[tnode] (n8) at (-1.6,-4.8) {$8$};
  \node[tnode] (n9) at ( 5.2,-4.8) {$9$};
  \node[tnode] (n10) at ( 5.2,-6.4) {$10$};
  \draw[tedge] (n1)--(n2) (n1)--(n3) (n2)--(n4) (n2)--(n5)
               (n3)--(n6) (n3)--(n7) (n5)--(n8) (n7)--(n9) (n9)--(n10);
  % cross-child pair (2,3): green arms
  \draw[armX] (n1) -- (n2);
  \draw[armX] (n1) -- (n3);
  % degenerate rows highlighted in orange (arm of length 0 sits at peak)
  \draw[arm0] (n1) -- (n2) -- (n4);
  \draw[arm0] (n2) -- (n5);
  \draw[arm0] (n1) -- (n3) -- (n6);
  \draw[arm0] (n3) -- (n7);
  \draw[arm0] (n2) -- (n5) -- (n8);
  \draw[arm0] (n3) -- (n7) -- (n9);
  \draw[arm0] (n7) -- (n9) -- (n10);
  \node[annot, font=\scriptsize, anchor=north west, text width=38mm, align=left]
    at (6.4,0.2)
      {\textbf{Peaks} (orange): $1,2,3,7$ carry $k{=}2$ rows.\\
       \textcolor{accentorange!70!black}{Orange paths}: arm-0 rides on $\ell$\\
       (degenerate, 7 rows).\\
       \textcolor{successgreen!60!black}{Green pair} at $\ell{=}1$: cross-child\\
       (strict, 3 rows). Total $=10$.};
\end{tikzpicture}
\caption{The $k=2$ census drawn on $T$. Every orange edge-chain
$\ell \to \cdots \to w$ with $\mathit{len}=2$ is a degenerate row (arm of
length~$0$ at $\ell$, endpoint $\ell$ itself); every green two-edge
splay $\{c_1, c_2\}$ at $\ell$ is a strict cross-child row. Peaks
$4,5,6,8,9,10$ are unmarked -- no $k=2$ pair files there. This figure is
the one Ch.16's \exref{c11-census} cites by name when it flags the
ancestor-of case.}\label{fig:c10-distk}
\end{figure}

\begin{tracebox}[$k=2$ census via \coderef{c10-distk-hist} on $T$, peak histogram path]
Post-order for the DFS of \coderef{c10-distk-hist} on $T$: 4, 8, 5, 2,
10, 9, 7, 6, 3, 1. Truncate histograms to length~$K+1=3$ (indices 0,1,2).
\\[2pt]
\textbf{Leaves}: $H[4]=(1,0,0)$. $H[6]=(1,0,0)$. $H[8]=(1,0,0)$.
$H[10]=(1,0,0)$.\\[2pt]
\textbf{Node 5} (child 8): harvest $H[5][0]\!\cdot\!H[8][1]=1\!\cdot\!0=0$.
Merge $H[5]=(1,1,0)$.\\
\textbf{Node 9} (child 10): harvest $0$ (same shape). $H[9]=(1,1,0)$.\\[2pt]
\textbf{Node 2} (children 4, 5): $H[2][0]=1$. After 4: harvest
$H[2][0]\!\cdot\!H[4][1]=1\!\cdot\!0=0$; merge $H[2]=(1,1,0)$. After 5:
harvest $H[2][0]\!\cdot\!H[5][1]+H[2][1]\!\cdot\!H[5][0]=1\!\cdot\!1+1
\!\cdot\!1=2$ (rows $\{2,8\}$ and $\{4,5\}$). Merge $H[2]=(1,2,1)$.\\[2pt]
\textbf{Node 7} (child 9): $H[7][0]=1$. Harvest $H[7][0]\!\cdot\!H[9][1]
= 1\!\cdot\!1 = 1$ (row $\{7,10\}$). Merge $H[7]=(1,2,0)$.\\
\textbf{Node 3} (children 6, 7): $H[3][0]=1$. After 6: harvest
$H[3][0]\!\cdot\!H[6][1]=0$. Merge $H[3]=(1,1,0)$. After 7: harvest
$H[3][0]\!\cdot\!H[7][1]+H[3][1]\!\cdot\!H[7][0]=1\!\cdot\!2+1\!\cdot\!1=3$.
\emph{Wait} -- that would put rows $\{3,9\}, \{6,7\}$ plus a phantom
$\{6,\!\cdot\!\}$; re-check: $H[7][1]=2$ counts nodes at distance 1 from
$7$ in $7$'s subtree, which are $\{9\}$ \emph{and} any other child --
here only $9$, so $H[7][1]=1$, and $H[7][0]=1$. Correction: $H[7]=(1,1,0)
\to$ merge yields $H[7]=(1,2,0)$ only \emph{after} merging 9, so at
harvest time $H[7]=(1,1,0)$. Recompute: $1\!\cdot\!1+1\!\cdot\!1=2$
(rows $\{3,9\}$ and $\{6,7\}$). Merge $H[3]=(1,2,1)$.\\[2pt]
\textbf{Node 1} (children 2, 3): $H[1][0]=1$. After 2: harvest
$H[1][0]\!\cdot\!H[2][1]=1\!\cdot\!2=2$ (rows $\{1,4\}, \{1,5\}$). Merge
$H[1]=(1,1,2)$. After 3: harvest $H[1][0]\!\cdot\!H[3][1]+H[1][1]
\!\cdot\!H[3][0]=1\!\cdot\!2 + 1\!\cdot\!1 = 3$ (rows $\{1,6\}, \{1,7\},
\{2,3\}$).\\[2pt]
\textbf{Total}: $2+2+1+2+2+3 = 10$. \checkmark\ matches
\ref{tab:c10-k2-table} and \figref{c10-distk} row-by-row.
\end{tracebox}

\begin{insightbox}
Seven of the ten rows in \ref{tab:c10-k2-table} are \emph{degenerate}:
one endpoint is the peak, one arm has length~$0$. The three strict
cross-child rows are the ones the naive $\sum_{c_1<c_2}\mathit{sz}(c_1)
\mathit{sz}(c_2)$ census catches; the seven degenerate rows are the ones
that vanish if $\mathit{H}[u][0]=1$ is not seeded before the child loop
(Trap~1 of \secref{c10-harvest}). Chapter~16's \exref{c11-census} cites
this exact distinction: $\lca(u,v)=u$ when $u\preceq v$ (the degenerate
case) and $\lca$ is a strict common ancestor otherwise (the cross-child
case). Both live in the same histogram; both are the same theorem.
\end{insightbox}

\subsection{Variations and Twists}

\begin{itemize}[nosep]
  \item \textbf{Range distance} $a\le d(u,w)\le b$: keep the peak
        histogram but sum over $\{i: a\le i+j+1\le b\}$. This is CSES
        2081 ``Fixed-Length Paths II'' (Ch.12 owns the URL). \hfill\lv{4}
  \item \textbf{Sum of distances to descendants at depth $d$}: replace
        the $\oplus$ in Shape~B by $+$ on integers -- the inverse is
        subtraction, so the backtracking still applies. \hfill\lv{3}
  \item \textbf{Non-invertible combiner} ($\max$, ``count of maximum''):
        Shape~B fails; fall back to Shape~A with small-to-large merges
        (Ch.17's \chref{c28}). \hfill\lv{4}
  \item \textbf{Unbounded $k$}: if the constraint is $\sum |H_v| \gg nk$,
        the histogram merge itself becomes the bottleneck --
        \chref{c27}'s knapsack-over-children lemma gives the amortised
        $\bigO{n\cdot\text{(budget)}}$ bound; \chref{c18}'s centroid
        recursion answers distance queries in $\bigO{n\log n}$.
        \hfill\lv{5}
\end{itemize}

\subsection{Common Pitfalls}

\begin{warningbox}[Trap 1 -- depth vs. distance (Ch.570D off-by-one)]
CF~570D's statement uses the convention ``the root has depth $1$'', and
``at depth $h$'' is \emph{absolute} (from the root), not relative to
$v$'s subtree. Writing \lstinline{depth[root]=0} in the code and
answering a query with \lstinline{X[h]} then reads one level too high.
Symptom: sample query \lstinline{(1,1)} returns ``Yes'' by accident,
deeper queries return ``No'' on hidden. Set the convention explicitly
and stick to it throughout the file.
\end{warningbox}

\begin{warningbox}[Trap 2 -- answering at entry instead of exit]
Shape~B must answer queries at $u$'s \emph{exit}, after every descendant
has been folded into $\mathit{X}[\cdot]$. Answering at entry reads the
accumulator before the subtree is processed, so queries see only the
ancestors of $u$ at that depth -- an entirely different (and wrong)
answer. Symptom: root-depth queries return correct values (the root's
``subtree at depth 1'' is just the root itself at exit-time), all other
queries come out too small.
\end{warningbox}

\begin{warningbox}[Trap 3 -- forgetting the shift in Shape~A's merge]
The line \lstinline{H[u][j] += H[v][j-1]} shifts by $+1$; the naive
\lstinline{H[u][j] += H[v][j]} counts each node at distance $j$ from
$v$ as being at distance $j$ from $u$. Symptom: the answer is right for
$k=1$ (shift of $+1$ doesn't matter at $k=1$ because $H[u][0]=1$ still
captures the ``$u$ paired with a child'' rows), wrong for every $k\ge 2$.
\end{warningbox}

\begin{warningbox}[Trap 4 -- \lstinline{int} popcount for wide masks]
Shape~B's XOR accumulator must be at least $2^C$-bit where $C$ is the
alphabet size. For lowercase letters ($C=26$), \lstinline{int} suffices
(32-bit); for arbitrary per-node weights as bit flags, promote to
\lstinline{long long}. Symptom: silent truncation of the character
bits, popcount comes out wrong only for letters past the 31st.
\end{warningbox}

\subsection{Problem Lab}

The section's anchor is CF~570D, \probref{c10-distk-solved}: $m$ offline
queries of the form ``at depth $h$ in $v$'s subtree, can the letters
form a palindrome?'' It is Shape~B straight off the page. Section~13.3
already solved the Shape~A instance (CF~161D); this lab rounds out the
pair so both distance-indexed shapes are on the record.

% ---------------------------------------------------------------------------
\solvedproblem{c10-distk-solved}{Tree Requests}{Codeforces 570D}
  {https://codeforces.com/problemset/problem/570/D}{CF 1800 L4}
  {distance-indexed DP, backtracking DFS, per-depth XOR accumulator, palindrome parity}

\begin{problemstatement}[Tree Requests (samples verified against the live statement)]
Roman planted a rooted tree of $n$ vertices. Vertex~$1$ is the root;
each vertex $i\ge 2$ has a parent $p_i < i$. Each vertex carries a
lowercase English letter. Depth of a vertex is the number of nodes on
the path from root to it, with the root at depth~$1$. Vertex~$u$ is in
the \emph{subtree} of $v$ if we can get from $u$ to $v$ by repeatedly
moving to the parent (in particular $v$ is in its own subtree). Given
$m$ queries $(v_i, h_i)$: consider the letters written on all vertices
in the subtree of $v_i$ located at depth $h_i$; can they be rearranged
into a palindrome? Constraints: $1\le n, m\le 5\cdot 10^{5}$; $1\le
v_i, h_i\le n$. Output: \texttt{Yes} or \texttt{No}, one line per
query. \textbf{Sample}: $n=6$, $m=5$; parents $1,1,1,3,3$ (for
vertices $2..6$); string \texttt{zacccd}; queries
$(1,1),(3,3),(4,1),(6,1),(1,2)$ $\Rightarrow$ \texttt{Yes No Yes Yes
Yes}.
\end{problemstatement}

\begin{attemptgate}[Attempt it yourself first -- \lv{4}, give yourself 20 minutes]
The naive per-query BFS costs $\bigO{n}$ each and dies at $n\cdot m =
2.5\cdot 10^{11}$. Two facts unlock the linear solution: (i) palindrome
$\iff$ at most one letter has odd count, i.e.\ popcount of a XOR mask
$\le 1$; (ii) queries are offline and grouped by $v$. Ask: what state
indexed by depth does a DFS maintain naturally along the way down?
\peek{Which property of a palindrome can you encode in a single 26-bit
      word?}
\nudge{The parity vector: bit $i$ of the word says whether letter~$i$
      has appeared an odd number of times. XOR accumulates parities
      for free, and popcount$\le 1$ is the test.}
\fullhint{Root the tree at~$1$ with depth starting at~$1$. Bucket the
      queries into \lstinline{q[v]} (each entry $(h, \mathit{id})$).
      Maintain a global \lstinline{X[d]} that is XOR-updated on entry
      to each node and \emph{never} restored on exit. On entering $v$,
      snapshot $\mathit{X}[h]$ for every query $(v,h)$ (this is the
      pre-subtree prefix); then XOR $v$'s own bit onto
      $\mathit{X}[\depth(v)]$ and recurse into children. On exiting
      $v$, answer each query as
      $\mathit{X}[h]\oplus\mathit{X}[h]_{\text{snapshot}}$ -- the
      difference is exactly the XOR of every $v$-subtree node at
      depth~$h$. Total $\bigO{n+m}$, memory $\bigO{n+m}$. See
      \thmref{c10-depth-xor} for why no-restore is the right shape.}
\end{attemptgate}

\solstep{Understanding the Problem}
Each query is a per-node, per-depth filter followed by a parity test.
A direct per-query BFS hits $\bigO{n}$ per query. Bucket queries by
$v$, then a single DFS visits every node in $v$'s subtree between
entering $v$ and leaving $v$; a per-depth \emph{monotone} accumulator
(never restored) maintained along the walk makes the difference
$\mathit{X}[h]_{\mathrm{out}}\oplus\mathit{X}[h]_{\mathrm{in}}$ the
subtree aggregate at $\bigO{1}$ per query.

\solstep{Key Observations}
\begin{observation}
A multiset of letters can be rearranged into a palindrome iff at most
one letter has odd count. Encode as a 26-bit XOR mask; the test is
$\mathit{popcount}(\mathit{mask}) \le 1$.
\end{observation}
\begin{observation}
DFS visits every node in $v$'s subtree, and only those, strictly
between $v$'s entry and exit. A per-depth accumulator that is
XOR-updated on entry but \emph{never restored} turns the query answer
into $\mathit{X}[h]_{\mathrm{out}}\oplus\mathit{X}[h]_{\mathrm{in}}$ --
$\bigO{1}$ per query.
\end{observation}
\begin{observation}
XOR is its own inverse, so \thmref{c10-depth-xor}'s group-invertibility
clause is satisfied -- the exit-time restore is a single XOR.
\end{observation}

\solstep{Think It Through With Me}
\begin{thinkbox}
\deadend{For each query, BFS the subtree of $v$, collect chars at depth
      $h$, count parities.}
{$\bigO{n\cdot m}$ in the worst case; at $n, m\le 5\cdot 10^{5}$ the
budget is $2.5\cdot 10^{11}$ -- TLE by four orders of magnitude, no
constant saves it.}
\deadend{Precompute for every $v$ a full depth-to-mask table
      $\mathit{Tab}[v][d]$.}
{Table is $\bigO{n^2}$ entries ($2.5\cdot 10^{11}$ masks), memory blows
up. Even with a small-to-large DSU merge (Ch.17) the total cost is
$\bigO{n\log n}$, but you'd be building the machinery before you need
it; the backtracking DFS is shorter and faster.}
\deadend{Euler-tour the tree, sort nodes by depth, binary-search each
      query's subtree interval.}
{Works, and it's the standard reference solution, but it uses Ch.24's
tin/tout machinery (\chref{c18}) which this chapter does not assume.
The backtracking DFS is the same asymptotic cost with fewer tools.}
\aha{Never restore the accumulator. Snapshot $\mathit{X}[h]$ per query
     at $v$'s entry; at $v$'s exit the XOR difference against the
     snapshot is the subtree aggregate. One DFS, no preprocessing.}
\end{thinkbox}

\solstep{Solution}
Build $G$ from parent pointers; depth of root~$1$ is~$1$. Bucket each
query $(v,h)$ into \lstinline{q[v]} tagged with its index. Set
$\mathit{X}[d]=0$ for every $d$. Run the DFS of \coderef{c10-depth-hist}:
at entry to $v$, snapshot $\mathit{X}[h]$ for every query in
$\mathit{q}[v]$; XOR $v$'s own bit onto $\mathit{X}[\depth(v)]$;
recurse into children; at exit, answer each query as
$\mathit{X}[h]\oplus\mathit{X}[h]_{\text{snapshot}}$. Do \emph{not}
restore the accumulator. Print answers in input order.

\solstep{Proof of Correctness}
\begin{proofbox}[\coderef{c10-depth-hist} answers every CF~570D query correctly]
Fix query $(v, h)$. \emph{Claim}: on exiting $v$, the difference
$\mathit{X}[h]\oplus\mathit{X}[h]_{\text{snapshot at $v$'s entry}}$
equals $\bigoplus_{u \in \mathrm{subtree}(v),\,\depth(u)=h} \mathit{bit}(u)$,
where $\mathit{bit}(u)$ is the mask \texttt{1 << (s[u] - 'a')}.
Because $\mathit{X}[\cdot]$ is never restored, its value at any
moment is the XOR of $\mathit{bit}(u)$ over every node $u$ already
visited in preorder, bucketed by $\depth(u)$. The DFS visits every
node in $v$'s subtree (and only those) between $v$'s entry and $v$'s
exit, so the two snapshots differ exactly by the subtree aggregate.
Formally: $\mathit{X}[h]_{\mathrm{out}} = \mathit{X}[h]_{\mathrm{in}}
\oplus \bigoplus_{u\in\mathrm{subtree}(v),\,\depth(u)=h}\mathit{bit}(u)$,
and XOR is its own inverse. \emph{Palindrome criterion}: a multiset of
letters can be rearranged into a palindrome iff at most one letter has
odd count, iff the popcount of the parity XOR mask is $\le 1$ (empty
subtree gives popcount~$0$ $\Rightarrow$ Yes, matching the statement's
Note). Both facts together yield the answer. \qed
\end{proofbox}

\solstep{Implementation}
\begin{lstlisting}[caption={CF 570D via iterative two-phase DFS with a per-depth prefix XOR accumulator (never restored); \lstinline{int} is enough for 26-bit masks}, label={lst:c10-570d-soln}]
#include <bits/stdc++.h>
using namespace std;

int main(){
    ios::sync_with_stdio(false); cin.tie(nullptr);
    int n, m; if(!(cin >> n >> m)) return 0;
    vector<vector<int>> adj(n + 1);
    for (int i = 2; i <= n; ++i){ int p; cin >> p; adj[p].push_back(i); }
    string s; cin >> s; s = " " + s;                       // 1-index letters
    vector<vector<pair<int,int>>> q(n + 1);                // q[v] = list of (h, id)
    for (int i = 0; i < m; ++i){ int v, h; cin >> v >> h; q[v].push_back({h, i}); }

    vector<int> X(n + 2, 0);                               // X[d] = 26-bit XOR prefix mask at depth d
    vector<int> in_val(m, 0);                              // entry snapshot of X[h] per query
    vector<char> ans(m, 0);
    vector<int> depth(n + 1, 0); depth[1] = 1;             // CF convention: root at depth 1

    // Two-phase iterative DFS to survive chains at n = 5e5.
    // Stack items are (node, phase). phase 0 = about to enter, phase 1 = about to exit.
    vector<pair<int,int>> st; st.push_back({1, 0});
    while (!st.empty()){
        auto [u, ph] = st.back(); st.pop_back();
        if (ph == 0){
            st.push_back({u, 1});                          // schedule exit next
            for (auto& [h, id] : q[u]) in_val[id] = X[h]; // snapshot BEFORE u's subtree
            X[depth[u]] ^= (1 << (s[u] - 'a'));           // add u's own bit
            // Push children in reverse so preorder matches the natural order.
            for (auto it = adj[u].rbegin(); it != adj[u].rend(); ++it){
                depth[*it] = depth[u] + 1;
                st.push_back({*it, 0});
            }
        } else {
            // Exit u: X[h] has accumulated everything in u's subtree (no child restored).
            for (auto& [h, id] : q[u]){
                int mask = (h <= n) ? (X[h] ^ in_val[id]) : 0;
                ans[id] = (__builtin_popcount((unsigned)mask) <= 1);
            }
            // NO restore of X[depth[u]] -- u's bit must remain visible to u's parent.
        }
    }
    for (int i = 0; i < m; ++i) cout << (ans[i] ? "Yes" : "No") << "\n";
    return 0;
}
\end{lstlisting}

\begin{complexitybox}
Time $\bigO{n+m}$: one XOR update per node in phase~0 and one
popcount per query in phase~1. Memory $\bigO{n+m}$ for adjacency,
query buckets, accumulator, and snapshots. At $n, m\le 5\cdot 10^{5}$:
$\le 10^{6}$ ops on the hot loop, comfortable under CF's 2-second
limit. The iterative two-phase stack is required -- a chain input at
$n=5\cdot 10^{5}$ blows a recursive frame. Verdict: \ac.
\end{complexitybox}

\solstep{Trace}
\begin{tracebox}[CF 570D sample: $n{=}6$, parents $1,1,1,3,3$, string \texttt{zacccd}, $m{=}5$ queries]
Depths: $1\!\to\!1$, $2\!\to\!2$, $3\!\to\!2$, $4\!\to\!2$, $5\!\to\!3$,
$6\!\to\!3$. Bit values $\mathit{bit}(u)=\texttt{1 << (s[u] - 'a')}$:
$1{=}\texttt{z}\mapsto 33554432$, $2{=}\texttt{a}\mapsto 1$,
$3{=}\texttt{c}\mapsto 4$, $4{=}\texttt{c}\mapsto 4$, $5{=}\texttt{c}\mapsto 4$,
$6{=}\texttt{d}\mapsto 8$. Children: $1\!\to\!\{2,3,4\}$,
$3\!\to\!\{5,6\}$. Queries bucketed:
$\mathit{q}[1]=\{(1,\text{id}0),(2,\text{id}4)\}$,
$\mathit{q}[3]=\{(3,\text{id}1)\}$, $\mathit{q}[4]=\{(1,\text{id}2)\}$,
$\mathit{q}[6]=\{(1,\text{id}3)\}$. Initial $\mathit{X}[\cdot]=0$.
\\[3pt]
Preorder: $1, 2, 3, 5, 6, 4$. Postorder (exit): $2, 5, 6, 3, 4, 1$. Trace
below shows $\mathit{X}[h]$ values that matter, and the snapshot
$\mathit{X}_{\text{in}}[h]$ read per query at entry.
\\[3pt]
\textbf{Enter 1}: snapshot $\mathit{in}_{\text{id}0}=\mathit{X}[1]=0$;
$\mathit{in}_{\text{id}4}=\mathit{X}[2]=0$. Add $1$'s bit:
$\mathit{X}[1]=33554432$.\\
\textbf{Enter 2}: no queries. Add $2$'s bit: $\mathit{X}[2]=1$.\\
\textbf{Exit 2}: no queries, no restore. $\mathit{X}[2]=1$.\\
\textbf{Enter 3}: snapshot $\mathit{in}_{\text{id}1}=\mathit{X}[3]=0$
(depth~$3$ still untouched). Add $3$'s bit: $\mathit{X}[2]=1\oplus 4=5$.\\
\textbf{Enter 5}: no queries. Add $5$'s bit: $\mathit{X}[3]=4$.\\
\textbf{Exit 5}: no queries.\\
\textbf{Enter 6}: snapshot $\mathit{in}_{\text{id}3}=\mathit{X}[1]=33554432$
(the only depth-$1$ node visited so far is $1$ itself). Add $6$'s bit:
$\mathit{X}[3]=4\oplus 8=12$.\\
\textbf{Exit 6}: query id3 asks $h{=}1$.
$\mathit{mask}=\mathit{X}[1]\oplus\mathit{in}_{\text{id}3}=33554432\oplus 33554432=0$.
popcount $=0\le 1 \Rightarrow$ \texttt{Yes} (id3).\\
\textbf{Exit 3}: query id1 asks $h{=}3$.
$\mathit{mask}=\mathit{X}[3]\oplus\mathit{in}_{\text{id}1}=12\oplus 0=12$.
popcount $=2 > 1 \Rightarrow$ \texttt{No} (id1).\\
\textbf{Enter 4}: snapshot $\mathit{in}_{\text{id}2}=\mathit{X}[1]=33554432$.
Add $4$'s bit: $\mathit{X}[2]=5\oplus 4=1$.\\
\textbf{Exit 4}: query id2 asks $h{=}1$.
$\mathit{mask}=\mathit{X}[1]\oplus\mathit{in}_{\text{id}2}=33554432\oplus 33554432=0$.
popcount $=0 \Rightarrow$ \texttt{Yes} (id2).\\
\textbf{Exit 1}: query id0 asks $h{=}1$.
$\mathit{mask}=\mathit{X}[1]\oplus\mathit{in}_{\text{id}0}=33554432\oplus 0=33554432$.
popcount $=1 \Rightarrow$ \texttt{Yes} (id0). Query id4 asks $h{=}2$.
$\mathit{mask}=\mathit{X}[2]\oplus\mathit{in}_{\text{id}4}=1\oplus 0=1$.
popcount $=1 \Rightarrow$ \texttt{Yes} (id4).
\\[3pt]
\textbf{Answers in input order}: id0 \texttt{Yes}, id1 \texttt{No},
id2 \texttt{Yes}, id3 \texttt{Yes}, id4 \texttt{Yes}. \\
\textbf{Cross-check with a direct subtree enumeration}: subtree $1$ at
depth~$1$ = $\{1\}$ (letter `z`, popcount~1 $\Rightarrow$ Yes);
subtree $3$ at depth~$3$ = $\{5,6\}$ (c$\oplus$d, popcount~2 $\Rightarrow$ No);
subtree $4$ at depth~$1$ = $\emptyset$ (Yes); subtree $6$ at depth~$1$
= $\emptyset$ (Yes); subtree $1$ at depth~$2$ = $\{2,3,4\}$
(a$\oplus$c$\oplus$c, popcount~1 $\Rightarrow$ Yes).
\textbf{Sample matches official} \texttt{Yes No Yes Yes Yes}. \\checkmark
\end{tracebox}

\solstep{Complexity}
$\bigO{n+m}$ time and $\bigO{n+m}$ memory, as the
\coderef{c10-570d-soln} complexity box states. The accumulator
$\mathit{X}[\cdot]$ is $\bigO{n}$ (depth $\le n$); the per-query
snapshot $\mathit{in}_{\text{id}}$ is $\bigO{m}$. No histogram is built
per node, so the $\bigO{nk}$ Space-A cost is completely sidestepped.

\solstep{Pitfalls and Variants}
\begin{warningbox}
Three ways to lose this problem: (i) \emph{restoring} $\mathit{X}[\depth(u)]$
on $u$'s exit -- the difference $\mathit{X}[h]_{\mathrm{out}}\oplus\mathit{X}[h]_{\mathrm{in}}$
then collapses to $0$ for every internal query, and only the deepest-
leaf queries survive; (ii) \emph{forgetting} the entry snapshot and
reading only $\mathit{X}[h]_{\mathrm{out}}$ -- the accumulator then
includes every preorder node before $u$ at depth $h$, so ancestors
contaminate the answer; (iii) depth-index confusion (root at $0$ vs.
root at $1$). Variant worth doing after this: CSES~2103 ``Company
Queries II'' -- same per-node offline accumulation, one line of binary
lifting added on the ancestor walk.
\end{warningbox}

\subsection{Practice Problems}

\begin{stuckbox}
\begin{itemize}[nosep, leftmargin=1.4em]
  \item Are queries bucketed \emph{at $v$} and answered \emph{at $u$'s
        exit}, not answered at entry? Entry reads only ancestors.
  \item Is $\mathit{X}[h]_{\text{in}}$ snapshotted per query before
        entering $u$'s subtree (not shared with other queries at
        the same $u$)? Sharing is fine if you compute the snapshot
        once per node; it is a bug to share across nodes.
  \item Is the root's depth $1$ (CF convention) or $0$? Pick one and
        stick to it in the whole file. Sample \lstinline{(1,1)}
        discriminates.
\end{itemize}
\end{stuckbox}

\begin{practiceset}
  \practice{Directing the Edges}{Codeforces 1385E}
    {https://codeforces.com/problemset/problem/1385/E}{CF 1700 L3}
    {Shape~A with a scalar accumulator. Root the tree at~$1$; count the
    minimum number of edges that must be flipped so every edge points
    away from some chosen root. Compute $\mathit{down}[v]$ (already-
    oriented edges in $v$'s subtree) and $\mathit{up}[v]$ (edges on the
    $v$-to-root path) in one DFS with the same exit-time aggregation
    spirit as \coderef{c10-depth-hist}; the answer is
    $\min_v(\mathit{down}[v]+\mathit{up}[v])$. Same fold skeleton as
    \secref{c10-subtrees}; the difference is that the state is a scalar
    pair, not a histogram.}
\end{practiceset}

\paragraph{Back-cites (URLs live elsewhere in the book).}
\begin{itemize}[nosep,leftmargin=1.4em]
  \item \textbf{CSES 2081 ``Fixed-Length Paths II''}: Ch.12's practice slot
        owns the pointer. This is Shape~A with a range condition
        $a \le d(u,w) \le b$ replacing the equality; the histogram merge
        is unchanged, only the sum's bounds shift. Re-solve it now that
        you know \thmref{c10-pairtop-count} is the general form.
  \item \textbf{LC 2049 ``Count Nodes With the Highest Score''}
        (\chref{c04} owns the URL): the removal-score definition is
        exactly the Ch.12 preview of Ch.15's rerooting. Ch.13 does not
        re-derive it; the fold is \secref{c10-subtrees}'s child template
        with a two-coordinate tuple.
\end{itemize}

\begin{recallbox}[10-second refresher -- 13.4 Distance-Indexed DP]
Two shapes of ``histogram over depth'': \textbf{Shape~A} = bottom-up
merge per node, $\bigO{nk}$, non-invertible combinable (CF~161D via
\thmref{c10-pairtop-count}); \textbf{Shape~B} = backtracking-DFS
accumulator with per-depth prefix XOR, $\bigO{n+m}$, invertible
combiners only (CF~570D via \thmref{c10-depth-xor}). The $k=2$
census on $T$ (\exref{c10-k2}, \figref{c10-distk}) is the reference
example Ch.16's \exref{c11-census} cites for the ``degenerate row''
-- $\ell$ is one endpoint, one arm has length~$0$; forgetting the
seed $\mathit{H}[u][0]=1$ drops those rows and is the top WA cause.
Knapsack-over-children budget analysis is Ch.14's (\chref{c27});
small-to-large and Euler-tour alternatives are Ch.17's (\chref{c28})
and Ch.24's (\chref{c18}).
\selftest{
  \item State \defref{c10-dist-hist} in one sentence. Why is the
        index the depth and the value the fold?
  \item Give one complexity line each for Shape~A and Shape~B, and
        name the algebraic property (invertibility) that Shape~B
        requires.
  \item In \ref{tab:c10-k2-table}, which rows are degenerate and
        which are strict cross-child? How many of each?
  \item Why does answering a CF~570D query at $u$'s \emph{entry}
        instead of $u$'s \emph{exit} produce a wrong result even
        when $\Phi$ is XOR?
  \item Which Ch.16 forward contract does \exref{c10-k2} discharge?
}
\end{recallbox}
